\chapter{Letzte Streuung und Baryogenese}

\section{Weiter hinauszuschauen heißt, in die Vergangenheit zu blicken}

In verschiedenen Entfernungen sehen wir das Universum zu verschiedenen Zeiten. In unserer Nähe ist die Rückblickzeit nicht besonders wichtig. Das Universum vor einem Jahr oder vor hundert Jahren unterschied sich kaum vom heutigen. Doch weit genug entfernt sehen wir das Universum in einem jüngeren Zustand. Ein jüngeres Universum ist dichter, und ein dichteres Universum ist heißer.

In einer hinreichend großen Entfernung war das Universum ungefähr so heiß wie die Sonnenoberfläche. Für uns sieht es nicht so aus, weil die Photonen auf ihrem Weg hierher rotverschoben wurden. Ein Rotverschiebungsfaktor von etwa tausend bedeutet, dass die beobachtete Temperatur um denselben Faktor niedriger ist:
\[
T_{\mathrm{obs}} \sim \frac{T_{\mathrm{emit}}}{1+z},
\qquad
1+z \sim 10^3.
\]\footnote{Diese Darstellung ist eine kompakte redaktionelle Transkription des mündlichen Arguments zur Rotverschiebung bei 00:02:20--00:02:39.}

Dort wurde das Universum undurchsichtig. Es bestand hauptsächlich aus Wasserstoff und zu einem kleinen Teil aus Helium. Bei Temperaturen von mehreren tausend Grad ist Wasserstoff ionisiert. Jenseits dieser Oberfläche gibt es also keine neutralen Atome, sondern Ionen: vor allem Protonen und Elektronen, die sich frei bewegen. Auch die Photonen sind dort energiereicher, im Wesentlichen Röntgenstrahlung. Falls sich zufällig ein Atom bildet, wird es durch Stöße oder ein Photon rasch wieder ionisiert.

Ein ionisiertes Gas ist elektrisch leitfähig, und ein Leiter lässt Strahlung nicht ungehindert hindurch. Ladungen reagieren auf elektrische Felder, bewegen sich hin und her und emittieren erneut Strahlung. Jenseits dieser Entfernung ist das Universum für uns also undurchsichtig. Die Oberfläche der letzten Streuung ähnelt der sichtbaren Oberfläche der Sonne: Sie ist keine materielle Wand, sondern eine Stelle, an der das Medium ziemlich abrupt von undurchsichtig zu durchsichtig wird.

Die meisten Teilchen dort sind Photonen, und zwar um etwa den Faktor \(10^8\). Heute gilt dasselbe, obwohl die Photonen inzwischen so energiearm sind, dass sie keinen bedeutenden Beitrag zur heutigen Energiedichte leisten. Sobald die Temperatur unter die Ionisationstemperatur des Wasserstoffs fällt, wird das Universum durchsichtig, und wir können hindurchsehen.

Dort endet der direkte Blick. Als Nächstes stellt sich die Frage, wie weit uns die bekannte Atom-, Kern- und Teilchenphysik noch in die Vergangenheit zurückführt.

\section{Noch weiter zurück}

Unmittelbar innerhalb der Oberfläche der letzten Streuung sind neben Photonen vor allem Protonen und Elektronen vorhanden, außerdem einige Heliumkerne und wenige Deuteronen. Antiprotonen und Positronen gibt es praktisch keine. Doch weiter zurück wird die Temperatur schließlich so hoch, dass zwei Photonen zusammenstoßen und ein Elektron und ein Positron erzeugen können:
\[
\gamma + \gamma \leftrightarrow e^- + e^+.
\]

Die grundlegende Schwelle ist erreicht, wenn die Photonen mindestens die Ruheenergie des Paares bereitstellen:
\[
E \ge 2m_e c^2.
\]
Unterhalb dieser Schwelle reicht gewöhnliches Licht dafür nicht aus. Oberhalb davon können Elektron--Positron-Paare entstehen, und auch die umgekehrte Reaktion kann stattfinden. Ein heißer Topf mit Materie bei dieser Temperatur enthielte Teilchenreaktionen, die dort die Rolle chemischer Reaktionen in gewöhnlicher Materie übernehmen.

Knapp oberhalb der Schwelle gibt es immer noch mehr Photonen als Elektronen und Positronen. Die Ruhemasse der Elektronen ist energetisch kostspielig. Steigt die Temperatur jedoch weit über die Schwelle, wird die Ruheenergie der Elektronen gegenüber ihrer kinetischen Energie unbedeutender, und die Dichten werden vergleichbar. Dann haben Elektronen, Positronen und Photonen ungefähr dieselbe Dichte und stoßen wie wild miteinander zusammen.

\begin{figure}[ht]
\centering
\includegraphics[width=0.62\textwidth]{lecture_06_figure_02.png}
\caption{Skizze der Schwelle für Paarerzeugung im Vergleich der Massenskalen von Elektron und Proton.}
\lectureframe{00:18:47}
\end{figure}

Als Nächstes stellt sich die Frage, was mit den Protonen geschieht. Sie sind bereits vorhanden und zu diesem Zeitpunkt sehr dicht, aber die thermische Energie reicht noch nicht aus, um Proton--Antiproton-Paare in großer Zahl zu erzeugen. Der Grund ist einfach:
\[
m_p \approx 2000\,m_e,
\qquad
E \ge 2m_p c^2
\]
Das ist eine viel höhere Schwelle.\editorialnote{In der Abbildung ist die Beschriftung der niedrigeren Masse weniger deutlich zu erkennen als die der oberen \(2m_e c^2\); im mündlichen Vergleich geht es jedoch ausdrücklich um die Schwellen für Elektronen und Protonen.}

Das ist bereits viel heißer als das Innere eines Sterns, in dem Antiteilchen nicht in nennenswerter Zahl vorkommen.

Doch auch das ist nur eine vorläufige Beschreibung. Sobald die Temperatur hoch genug ist, um Protonen und Antiprotonen frei zu erzeugen, reicht sie auch aus, um Protonen in Quarks zu zerlegen. Die wirklich heiße Suppe lässt sich daher besser als eine Suppe aus Elektronen, Positronen, Quarks, Antiquarks, Photonen, Neutrinos, Gluonen und weiteren Teilchensorten beschreiben. Eine einfache schematische Reaktion von der Tafel lautet:
\[
e^- + e^+ \to \gamma \to q + \bar q.
\]

\begin{classroomqa}
\lecturetimestamp{00:20:05}
\audiencequestion{Sind Mesonen gemeint oder Vorquarks?}
\lecturerresponse{Mesonen bestehen aus Quarks. Ist die Temperatur hoch genug, um ein Proton zu zerlegen, zerlegt sie auch die Mesonen.}
\end{classroomqa}

Bis hierhin haben wir die Entwicklung rückwärts in der Zeit verfolgt. Drehen wir sie nun um und fragen, was geschieht, wenn die heiße symmetrische Suppe abkühlt.

\section{Elektrische Neutralität und ein symmetrisches heißes Universum}

Bei hinreichend hoher Temperatur enthält das Universum gewaltige Dichten von Teilchen und Antiteilchen. Elektronen und Positronen treten paarweise auf, Quarks und Antiquarks ebenfalls, und naheliegend ist die Annahme, dass jede Teilchensorte und ihr Antiteilchen nahezu gleich häufig vorkommen.

Bei der elektrischen Neutralität ist eine kleine Präzisierung nötig. Das Universum ist heute elektrisch neutral. Als Eselsbrücke kann man sich nicht unser tatsächliches Universum vorstellen, sondern eines, das geschlossen und endlich ist, etwa eine Dreisphäre. Betrachten wir zunächst als Modell eine zweidimensionale Kugeloberfläche. Elektrische Feldlinien können nur an Ladungen beginnen oder enden. Eine positive Ladung sendet Feldlinien aus; eine negative Ladung nimmt sie auf. Auf einer geschlossenen Kugeloberfläche können diese Linien nicht einfach ins Unendliche entweichen. Sie bündeln sich auf der gegenüberliegenden Seite wieder. Als Eselsbrücke gilt daher: Ein geschlossenes und endliches Universum kann keine Nettoladung tragen. Ob unser tatsächliches Universum buchstäblich geschlossen ist, ist hier nicht entscheidend. Entscheidend ist, dass sich positive und negative Ladungen auf großen Skalen offenbar außerordentlich genau aufheben.

Elektrische Neutralität bedeutet \emph{nicht}, dass es gleich viele Protonen und Antiprotonen gibt. Sie bedeutet, dass die Gesamtladung aller positiven Teilchen der Gesamtladung aller negativen Teilchen entspricht. Protonen, Positronen und alle weiteren positiv geladenen Teilchensorten gleichen gemeinsam die Ladung der Elektronen und aller weiteren negativ geladenen Teilchensorten aus.

Soweit wir der bekannten Physik vertrauen können, und vielleicht noch ein wenig darüber hinaus, war das Universum heiß genug, um Teilchen und Antiteilchen in großer Zahl hervorzubringen. Zu jedem Beschleunigerexperiment, das man sich vorstellen kann, gab es irgendwann im frühen Universum ein Gegenstück. Ein nahezu symmetrischer Anfang ist die naheliegende Vermutung: gleich viele Elektronen und Positronen, gleich viele Quarks und Antiquarks und, wenn Protonen und Antiprotonen als Stellvertreter für Baryonen und Antibaryonen dienen,
\[
N_p = N_{\bar p}.
\]

Diese Vermutung ist attraktiv, weil die physikalischen Gesetze zwischen Teilchen und Antiteilchen nahezu symmetrisch sind. Ein Positron verhält sich fast genauso wie ein Elektron, nicht ganz, aber sehr ähnlich. Genau diese nahezu bestehende Symmetrie ist das Problem.

\section{Den Film vorwärts abspielen}

Nehmen wir an, das Universum beginnt in einem heißen Zustand mit gleich vielen oder nahezu gleich vielen Teilchen und Antiteilchen und kühlt dann ab.

Solange die Temperatur hoch genug ist, können Protonen und Antiprotonen bei energiereichen Stößen entstehen und wieder zu Photonen annihilieren. Doch irgendwann fällt die Temperatur unter die Schwelle für Protonen. Dann können Proton--Antiproton-Paare weiterhin zu Photonen annihilieren, aber nach einer kurzen weiteren Expansion haben diese Photonen so viel Energie verloren, dass sie keine Proton--Antiproton-Paare mehr erzeugen können. Das ist eine Einbahnstraße. Später geschieht dasselbe mit Elektronen und Positronen, bei einer viel niedrigeren Temperatur als der Schwelle für Protonen, aber immer noch weit oberhalb von Sterntemperaturen.

Da die Expansion im Vergleich zu mikroskopischen Reaktionsraten langsam verläuft, bleibt das Universum währenddessen größtenteils nahe am thermischen Gleichgewicht. Das ist hier mit adiabatischer Expansion gemeint. Genauer gesagt ändert sich die Entropie dabei nur sehr wenig. In einer heißen Teilchensuppe entspricht die Entropie im Wesentlichen der Teilchenzahl, bis auf Faktoren der Größenordnung eins. Die annihilierenden Protonen, Antiprotonen, Elektronen und Positronen werden also ungefähr durch Photonen ersetzt.

Wären Protonen und Antiprotonen genau gleich häufig, würden sie vollständig annihilieren. Dasselbe geschähe mit Elektronen und Positronen. Übrig blieben Photonen. Gäbe es jedoch auch nur einen winzigen Überschuss an Protonen gegenüber Antiprotonen, würden die entsprechenden Paare annihilieren und der nicht gepaarte Überschuss bliebe bestehen. Ein konkretes Beispiel macht das deutlich: Gäbe es hundert Protonen und fünfzig Antiprotonen, würden die fünfzig Antiprotonen jeweils ein Proton finden und mit ihm annihilieren; fünfzig Protonen blieben übrig. Die Ladungsneutralität erfordert dann einen entsprechenden verbleibenden Überschuss an Elektronen gegenüber Positronen.

Aus der heutigen Photonenzahl lässt sich so auf die frühe Asymmetrie schließen. Heute ist die Zahl der Photonen etwa \(10^9\) bis \(10^{10}\) Mal so groß wie die Zahl der Protonen oder Elektronen:
\[
\frac{N_\gamma}{N_p} \sim 10^9\text{--}10^{10}.
\]
In der frühen heißen Suppe gab es also nur etwa ein überschüssiges Proton pro Milliarde Proton--Antiproton-Paare. Dieses winzige Ungleichgewicht ist der ganze Grund dafür, dass es überhaupt noch Atome gibt.

\begin{classroomqa}
\lecturetimestamp{00:42:25}
\audiencequestion{Warum sollte man den Überschuss nicht statistisch verstehen, als etwas, das beim Abkühlen mit einer winzigen Wahrscheinlichkeit entsteht?}
\lecturerresponse{Um aus einem anfangs symmetrischen Zustand ein zusätzliches Proton zu erhalten, müsste man ein Proton erzeugen, ohne zugleich ein Antiproton zu erzeugen. Außerdem müsste man damit eine negative Ladung erzeugen, etwa ein Elektron. Das bekannte Standardmodell enthält keinen solchen Prozess bei niedrigen Energien. Bei den untersuchten Energien bleibt die Baryonenzahl erhalten; gewöhnliches Abkühlen erzeugt daher keinen Nettobaryonenüberschuss.}
\end{classroomqa}

Die elektrische Ladung ist also kein Hindernis. Ein schematisches Beispiel wäre
\[
\gamma + \gamma \to p + e^-.
\]\footnote{Diese Darstellung ist eine redaktionelle Transkription des mündlichen Beispiels bei 00:44:16--00:46:35. Sie soll verdeutlichen, dass die elektrische Ladung erhalten bleiben könnte, während die Baryonenzahl verletzt wird.}
Insgesamt wäre dies elektrisch neutral. Es verletzt jedoch eine andere empirische Regel. Im untersuchten Standardmodell kann man kein Proton erzeugen, ohne zugleich ein Antiproton zu erzeugen. Für unsere Zwecke ist die relevante Erhaltungsgröße die Baryonenzahl:
\[
B=(N_p+N_n)-(N_{\bar p}+N_{\bar n}).
\]\footnote{Diese kompakte Formel präzisiert die mündliche Definition bei 00:45:18--00:45:40.}
Bei den bekannten Energien ist diese Größe erhalten. Die Asymmetrie kann daher während des späten, gut untersuchten Abschnitts der Abkühlung nicht entstehen.

\section{Baryogenese}

Darum geht es bei der sogenannten Baryogenese. Nehmen wir an, das Universum hätte tatsächlich mit gleich vielen Baryonen und Antibaryonen begonnen. Was müssten die physikalischen Gesetze zulassen, damit aus diesem symmetrischen Zustand der winzige beobachtete Überschuss entsteht?

An der Tafel wurde dies auf zwei symbolische Ausdrücke reduziert,
\[
N_p=N_{\bar p},
\qquad
N_p-N_{\bar p},
\]
und anschließend auf eine Liste von Bedingungen, die den ersten Ausdruck mit dem zweiten verknüpfen.

\begin{figure}[ht]
\centering
\includegraphics[width=0.72\textwidth]{lecture_06_figure_04.png}
\caption{Bedingungen für die Erzeugung einer Baryonenasymmetrie aus einem anfangs symmetrischen Zustand.}
\lectureframe{01:22:26}
\end{figure}

\editorialnote{In dieser Aufnahme sind die Wörter der nummerierten Liste nur teilweise lesbar. Die folgenden drei Punkte entsprechen der begleitenden mündlichen Erläuterung.}

Die erste Bedingung ist die Verletzung der Baryonenzahl. Ist die Baryonenzahl eine exakte Erhaltungsgröße, bleibt sie in einem Universum, das mit einer Nettobaryonenzahl von null beginnt, auch weiterhin null.

Die zweite Bedingung ist eine Bevorzugung von Teilchen gegenüber Antiteilchen oder umgekehrt. Eine Verletzung der Baryonenzahl allein genügt nicht. Lässt die Natur einen protonenerzeugenden Prozess und den entsprechenden spiegelbildlichen antiprotonenerzeugenden Prozess mit gleicher Wahrscheinlichkeit zu, ist das mittlere Ungleichgewicht weiterhin null. Die Reaktionsraten müssen sich unterscheiden.

Gewöhnliche statistische Fluktuationen sind ebenfalls viel zu klein. Ist die Gesamtzahl der Teilchen, einschließlich der Photonen, ungefähr \(10^{90}\), dann ergibt eine Fluktuation der Größenordnung \(\sqrt{N}\) nur etwa \(10^{45}\). Das reicht bei Weitem nicht aus, um die ungefähr \(10^{80}\) Protonen im Universum zu erklären.

Die dritte Bedingung lautet, dass sich das Universum irgendwann hinreichend weit vom Gleichgewicht entfernt haben muss. Hier stützt sich das Argument auf den CPT-Satz. Wenn ein heißes Gas schon immer im Gleichgewicht war, sieht ein rückwärts abgespielter Film genauso aus wie ein vorwärts abgespielter. In dieser Situation wird die Symmetrie zwischen Teilchen und Antiteilchen durch die exakte Symmetrie, die Ladungskonjugation, Parität und Zeitumkehr verbindet, effektiv wiederhergestellt. Damit eine tatsächliche Asymmetrie bestehen bleibt, muss der reale Verlauf auf mikroskopischen Zeitskalen vorwärts und rückwärts unterschiedlich aussehen. Zu einem ausreichend frühen Zeitpunkt war die Expansion des Universums schnell genug, um das zu bewirken.

Die drei Bedingungen lauten also:
\begin{enumerate}
\item Verletzung der Baryonenzahl;
\item Verletzung der Symmetrie zwischen Teilchen und Antiteilchen, genauer gesagt CP-Verletzung;
\item eine Phase außerhalb des thermischen Gleichgewichts.
\end{enumerate}

\begin{classroomqa}
\lecturetimestamp{01:09:08}
\audiencequestion{Ordnet man diese Bedingungen CPT zu: Bezieht sich Bedingung 2 tatsächlich auf CP und Bedingung 3 auf die Zeitrichtung?}
\lecturerresponse{Ja. Bedingung 2 sollte eigentlich als CP-Verletzung formuliert werden. Bedingung 3 bezeichnet den Umstand, dass der tatsächliche Verlauf außerhalb des Gleichgewichts lag; sie ist weder ein eigenes Erhaltungsgesetz noch dasselbe wie eine Verletzung der Baryonenzahl.}
\end{classroomqa}

Sind diese drei Bedingungen erfüllt, kann auch ein Universum, das mit gleich vielen Protonen und Antiprotonen beginnt, in die eine oder andere Richtung verschoben werden. Ungewiss bleibt, ob die Natur die Baryonenzahl tatsächlich verletzt und, falls ja, auf welche Weise.

\section{Was bekannt ist und was nicht}

Der offene Teil der Geschichte betrifft die Verletzung der Baryonenzahl. Wir haben noch nie in der Natur eine Reaktion direkt beobachtet, bei der sich
\[
(N_p+N_n)-(N_{\bar p}+N_{\bar n}).
\]
ändert. Das könnte lediglich bedeuten, dass die betreffenden Reaktionen viel höhere Temperaturen erfordern, als im Labor direkt untersucht werden können, oder dass sie bei niedrigen Energien äußerst unwahrscheinlich sind. Bislang bleibt dies jedoch die offene Frage.

Darum ist der Protonenzerfall in dieser Diskussion so wichtig. Ist die Baryonenzahl nicht exakt erhalten, müsste das Proton prinzipiell instabil sein. Tatsächlich ist es aber außerordentlich stabil. Die hier angegebene Untergrenze liegt in der Größenordnung von
\[
\tau_p \gtrsim 10^{34}\ \text{years}.
\]
Sie stammt aus riesigen Experimenten, bei denen sehr große Wasservolumina mit Lichtdetektoren ausgestattet und tief unter der Erde platziert wurden. Sie widerlegt eine Verletzung der Baryonenzahl nicht. Sie bedeutet lediglich, dass jede solche Verletzung bei niedrigen Energien außerordentlich stark unterdrückt sein muss.

In der Diskussion kommt auch eine weitere Bemerkung zur bekannten Physik vor. Bei etwa \(100\) bis \(250\,\mathrm{GeV}\) ähneln die schwachen Wechselwirkungen den elektromagnetischen Wechselwirkungen deutlich stärker. Das könnte in detaillierten Modellen eine Rolle spielen, ist aber nicht der entscheidende Schritt im vorliegenden Argument.

Dunkle Materie ändert an der grundlegenden Logik nicht viel. In den meisten Theorien wechselwirkt sie so schwach mit gewöhnlicher Materie, dass sie für diese Frage praktisch nur eine Nebenrolle spielt.

\begin{classroomqa}
\lecturetimestamp{01:24:49}
\audiencequestion{Wie beeinflusst die Gravitation all das?}
\lecturerresponse{Vor allem auf zweierlei Weise. Erstens gäbe es ohne die Expansion des Universums keine Phase außerhalb des Gleichgewichts. Zweitens liefert die Verdampfung Schwarzer Löcher einen starken Grund für die Vermutung, dass die Baryonenzahl nicht für alle Zeiten eine exakte Erhaltungsgröße sein kann. Davon abgesehen sind gewöhnliche Gravitationswechselwirkungen zwischen Teilchen hier zu schwach, um eine große Rolle zu spielen.}
\end{classroomqa}

Anschließend stellt sich eine weitere Frage zur elektrischen Ladung. Würde dasselbe Argument mit Schwarzen Löchern auch die Ladungserhaltung infrage stellen? Nein. Ein geladenes Schwarzes Loch trägt elektrische Flusslinien, die nach außen reichen und die Ladung kenntlich machen. Ein solches Schwarzes Loch verdampft nicht vollständig; stattdessen nähert es sich einem extremalen Überrest. Das Argument mit Schwarzen Löchern stellt also die Baryonenzahl infrage, nicht die elektrische Ladung.

Eine weitere Frage lautet, ob man die sehr dichte frühe Quarksuppe einfach als Schwarzes Loch auffassen sollte. Nein. Es geht nicht um einen statischen, dichten Klumpen. Das Universum expandierte. Die Geschichte mit Schwarzen Löchern ist ein Gedankenexperiment zur Frage, was aus der Baryonenzahl wird, und keine wörtliche Beschreibung des frühen Universums als eingeschlossenes Gebiet.

\begin{classroomqa}
\lecturetimestamp{01:31:14}
\audiencequestion{Wüsste man, wie man das Standardmodell um solche Reaktionen ergänzen könnte?}
\lecturerresponse{Das Problem ist eher das Gegenteil. Man kann leicht zu starke Verletzungen der Baryonenzahl hinzufügen und dadurch den Protonenzerfall viel zu schnell machen. Die eigentliche Schwierigkeit besteht darin, das Proton so stabil zu halten, wie es tatsächlich ist, und zugleich die gewünschten Prozesse im frühen Universum zu ermöglichen.}
\end{classroomqa}

\section{Entropie und die späteren Fragen zu Neutronen}

Kehren wir nun zur Bilanz zurück. In einem adiabatisch expandierenden Universum bleibt die Entropie annähernd erhalten. In einer heißen Suppe aus Elementarteilchen entspricht die Entropie im Wesentlichen der Teilchenzahl, bis auf Faktoren der Größenordnung eins. Deshalb muss man keine gewaltige späte Erzeugung von Photonen aus dem Nichts annehmen. Bei den Annihilationen wird der Inhalt der Suppe größtenteils zu Strahlung umverteilt.

In der älteren Fassung des Problems war die Geschichte auf den Kopf gestellt. Statt zu fragen, warum im Verhältnis zu den Photonen so wenige Protonen überlebt haben, fragte man umgekehrt, warum es im Vergleich zu den Protonen so viele Photonen gibt. Die präzisere Frage lautet jedoch nicht, warum das Universum spät im Verlauf so viel Entropie erzeugt hat. Vielmehr geht es darum, warum die Asymmetrie zwischen Teilchen und Antiteilchen von Anfang an so winzig war: nur etwa ein überschüssiges Proton pro Milliarde Proton--Antiproton-Paare.

Ein freies Neutron zerfällt zwar, doch dieser Prozess \emph{verletzt nicht} die Baryonenzahl:
\[
n \to p + e^- + \bar\nu.
\]
Das Neutron trägt dieselbe Baryonenzahl wie das Proton.

Bei einem freien Neutron ist die verfügbare Zerfallsenergie sehr klein, weil die Neutronenmasse nur geringfügig über der Masse des Endzustands liegt. Das ist einer der Gründe, weshalb der Zerfall gemessen an den Maßstäben der Elementarteilchenphysik so langsam verläuft. Ist ein Neutron in einem Kern gebunden, kann die Bindungsenergie diese geringe verfügbare Energie vollständig aufzehren; dann kann das gebundene Neutron stabil sein.

\begin{classroomqa}
\lecturetimestamp{01:43:03}
\audiencequestion{Verschwinden Neutronen einfach aus dem Universum, oder entstehen auch neue?}
\lecturerresponse{Zunächst lautet die Antwort, dass freie Neutronen durch Zerfall verloren gehen. Das ist aber nicht die ganze Geschichte. Es gibt tatsächliche Quellen für die Entstehung von Neutronen, darunter Kernreaktionen in Sternen und die Bildung von Neutronensternen. Die Gesamtzahl der Neutronen nimmt also nicht einfach fortwährend ab.}
\end{classroomqa}